\begin{abstract}
  Hofmann (1999) introduced the functional programming language \lfpl to characterize the functions computable in polynomial time using an affine type system.
  %
  \lfpl enables a natural programming style, including nested recursion, and has inspired the development of type systems for automatic cost analysis, linear dependent type theories, and efficient memory management in functional programming languages.
  %
  Despite its prominence, there does not exist a self-contained presentation, let alone a full mechanization, of \lfpl and its core metatheory.
  
  This article presents a modern account and mechanization of \lfpl and its metatheory with the goal of being self-contained and accessible while presenting the strongest-known soundness and completeness results.
  %
  The soundness proof works with the language \lfplplus, which extends \lfpl with additional language features, thereby providing a stronger theorem than pre-existing results.
  %
  The proof is novel, adapting a technique by Aehlig and Schwichtenberg (2002) to construct explicit polynomials that bound the cost of an \lfplplus expression with respect to a big-step cost semantics.
  %
  The completeness proof shows that \lfpl programs can simulate polynomial-time Turing machines while only relying on restricted forms of linear functions and lists.
  %
  It has the same structure as the original proof by Hofmann (2002) but greatly simplifies the core argument with a novel stack-like data structure, implemented with first-class functions and lists.
  %
  The mechanization includes the full soundness and completeness proofs and serves as one of the first case studies of mechanized metatheory in the recently-developed proof assistant Istari.
  %
  Takeaway?
\end{abstract}
